\problemname{Hair Dyeing}

Your friend has grown tired of their natural hair color and has decided it's time for a change.

Every day, your local hair salon has a big discount on one of its $10$ hair dyes (none of which
match your friend's natural hair color). They maintain a preliminary schedule of their discounts
for the next $N$ days but warn that the schedule may change in the future. Dyeing hair is, of
course, not something to be done half-heartedly, so your friend wants to make a plan for how to
dye their hair in the coming days.

Such a plan determines on which days your friend will go to dye their hair. When dyeing, they
always choose the discounted color since it would be too expensive otherwise. In total, there
are $2^N$ possible plans. However, they do not accept all plans: naturally, they do not want
to dye their hair the same color two days in a row.

You have been asked by your friend to calculate how many different such plans they could create.
However, each time you provide an answer, your friend informs you that the salon has updated its
schedule, changing which color is on discount on one of the days.

Since the answer can be very large, you should output it modulo $2^{64}$. This means the remainder
of the answer after dividing by $2^{64}$. In most programming languages, all calculations are
automatically performed modulo $2^{64}$ when using unsigned 64-bit integers.
In C++, you can use \texttt{unsigned long long}.

\section*{Input}
The first line of input contains two integers, $N$ and $Q$ ($1 \leq N \leq 10^5$, $0 \leq Q \leq 10^5$),
representing the number of days in the schedule and the number of schedule updates.

Next, there are $N$ integers $a_i$ ($0 \leq a_i \leq 9$), where $a_i$ indicates that color $a_i$ is on
discount on day $i$.

After that, there are $Q$ lines, each containing the integers $i$ and $v$ ($1 \leq i \leq N$, $0 \leq v \leq 9$),
meaning that on day $i$, the discount is now on color $v$.

\textbf{Note} that both $a_i$ and $v$ are very small.

\section*{Output}
First, print the number of hair-dyeing plans that can be created from the original schedule, modulo $2^{64}$.

Then, after each update, print the number of plans that can be created modulo $2^{64}$.

\section*{Scoring}
Your solution will be tested on a set of test groups, each worth a number of points. Each test group contains
a set of test cases. To get the points for a test group you need to solve all test cases in the test group.

\noindent
\begin{tabular}{| l | l | p{12cm} |}
  \hline
  \textbf{Group} & \textbf{Points} & \textbf{Constraints} \\ \hline
  $1$    & $3$        & $Q=0, N \leq 15$ \\ \hline
  $2$    & $7$        & $Q=0, N \leq 1000$ \\ \hline
  $3$    & $8$        & $Q=0$ \\ \hline
  $4$    & $15$       & $Q \leq 6000$ \\ \hline
  $5$    & $30$       & $a_i, v \leq 1$ \\ \hline
  $6$    & $20$       & $a_i, v \leq 3$ \\ \hline
  $7$    & $7$        & $a_i, v \leq 6$ \\ \hline
  $8$    & $10$       & No additional constraints. \\ \hline
\end{tabular}


\section*{Explanation of Sample 1}
There are 4 plans: either the friend can choose not to dye their hair at all or dye their hair exactly once.

\section*{Explanation of Sample 2}
The 7 plans are:
\begin{itemize}
  \item $[]$
  \item $[0]$
  \item $[1]$
  \item $[0 1]$
  \item $[0]$
  \item $[1 0]$
  \item $[0 1 0]$
\end{itemize}
Note that plans that contain only the color $0$ are counted twice since dyeing hair on different days counts
as different plans, even if the color is the same.

